\documentclass{beamer}

\usepackage{fontspec}
\usepackage[ngerman]{babel}
\usepackage{csquotes}

\usetheme{FSFW}

\begin{document}

\title{Digitale Selbstverteidigung}
\subtitle{Freie Software und Freie Inhalte für die Schule}

\maketitle

\begin{frame}
  \frametitle{Software in der Schule}

  \onslide<+->

  \begin{block}{Problem}
    Verwendung von Software in der Schule wünschenswert, aber
    nicht-technische Aspekte schwierig
    \begin{itemize}
    \item Kosten
    \item Lizenzierung
    \item Datenschutz
    \item Gleichberechtigung
    \end{itemize}
  \end{block}

  \onslide<+->

  \begin{block}{Achtung!}
    Aspekt \enquote{kostenlos} allein nicht ausreichend (mögliche
    Lizenzprobleme)
  \end{block}

  \onslide<+->

  \begin{block}{Alternativen}
    \emph{Freie Software} statt proprietär Programme
    \begin{itemize}
    \item meist nicht so \enquote{gut} wie kommerzielle Programme
    \item für den Schulgebrauch aber meist mehr als ausreichend!
    \end{itemize}
  \end{block}

\end{frame}

\begin{frame}
  \frametitle{Beispiel~1: LibreOffice}

  \onslide<+->

  \begin{block}{Was'n das?}
    \begin{itemize}
    \item freie Implementierung einer Office Suite (Textverarbeitung,
      Tabellenkalkulation, Präsentationen, \dots)
    \item verfügbar für Windows, Mac, und Linux
    \end{itemize}
  \end{block}

  \onslide<+->

  \begin{block}{Warum?}
    \begin{itemize}
    \item für alle Schüler ohne Probleme verfügbar
    \item Unterstützung offener Dateiformate, die auch durch andere Programme
      lesbar sind
    \item keine Abhängigkeit von einem Hersteller
    \item keine Lizenzierung notwendig (insbesondere kostenlos)
    \end{itemize}
  \end{block}

  \onslide<+->

  \begin{block}{Woher?}
    \url{https://www.libreoffice.org/}
  \end{block}

  % TODO: screenshots from https://www.libreoffice.org/discover/screenshots/

\end{frame}

\begin{frame}
  \frametitle{Beispiel~2: OwnCloud/NextCloud}

  \onslide<+->

  \begin{block}{Was'n das?}
    \begin{itemize}
    \item Freies Programm zum Datenaustausch und zur Datensynchronisation
    \item verfügbar für Windows, Mac, Linux, und Mobilgeräte
    \end{itemize}
  \end{block}

  \onslide<+->

  \begin{block}{Warum?}
    \begin{itemize}
    \item Kontrolle über eigene Daten
    \item Einhaltung geltender Datenschutzregelungen
    \item keine Lizenzierung notwendig (insbesondere kostenlos)
    \end{itemize}
  \end{block}

  \onslide<+->

  \begin{block}{Woher?}
    \begin{itemize}
    \item \url{https://nextcloud.com/}
    \item \url{https://owncloud.org/}
    \end{itemize}
  \end{block}

  % TODO: screenshots

\end{frame}

\begin{frame}
  \frametitle{Beispiel~3: Gimp und Inkscape}

  \onslide<+->

  \begin{block}{Was'n das?}
    \begin{itemize}
    \item Bildbearbeitung; pixelbasiert (Gimp) oder vektorbasiert (Inkscape)
    \item erhältlich für Windows, Mac, und Linux
    \end{itemize}
  \end{block}

  \onslide<+->

  \begin{block}{Warum?}
    \begin{itemize}
    \item für Schulgebrauch mehr als ausreichend
    \item für alle Schüler ohne Probleme verfügbar
    \item keine Lizenzierung notwendig (insbesondere kostenlos)
    \end{itemize}
  \end{block}

  \onslide<+->

  \begin{block}{Woher?}
    \begin{itemize}
    \item \url{https://gimp.org}
    \item \url{https://inkscape.org}
    \end{itemize}
  \end{block}

  % TODO: screenshots

\end{frame}

\begin{frame}
  \frametitle{Weitere Beispiele}

  \onslide<+->

  \begin{itemize}
  \item GNU/Linux, ein freies Betriebssystem % eventuell skolelinux nennen?
                                % scheint aber nicht ganz aktuell zu sein :(
  \item GeoGebra (\url{https://www.geogebra.org/})
  \item Stellarium (\url{www.stellarium.org})
  \item Minetest, eine freie Alternative zu Minecraft (\url{minetest.net})
  \item Jabber/XMPP, ein freies und dezentrales Instant Messaging Protokoll
  \item Scribus, ein freies Publishing-Programm (\url{https://www.scribus.net/})
  \item ShotCut/OpenShot für Videobearbeitung (\url{https://www.shotcut.org/})
  \item Audacity für Audiobearbeitung (\url{www.audacityteam.org/})
  \item diaspora*, ein soziales Netzwerk (\url{https://www.joindiaspora.com/})
  \item GNUSocial, Quitter, Mastodon, \dots{} Kurznachrichtendienste (\url{https://mastodon.social})
  \item Notepad++, ein freier Texteditor (\url{https://notepad-plus-plus.org/})
  \end{itemize}

\end{frame}

\begin{frame}
  \frametitle{Software-Freiheiten}

  \onslide<+->
  \onslide<+->

  \begin{block}{Software für jeden Anforderung benutzen zu können}
    \begin{itemize}
    \item Erlaubt die uneingeschränkte Verwendung von Programmen
    \item Schützt vor komplexen Nutzungsvereinbarungen und versehentlichen
      Lizenzverstößen; Verwendung wird erheblich vereinfacht.
    \end{itemize}
  \end{block}

  \onslide<+->

  \begin{block}{Software untersuchen und modifizieren zu können}
    \begin{itemize}
    \item Erlaubt das Verständnis vorhandener Programme, die Anpassung an
      spezielle Anforderungen, und die Überprüfung der Korrektheit
    \item Schützt vor Abhängigkeiten und unerwünschtem Verhalten
    \end{itemize}
  \end{block}

\end{frame}

\begin{frame}
  \frametitle{Software-Freiheiten}

  \onslide<+->

  \begin{block}{Kopien der Software weitergeben zu können}
    \begin{itemize}
    \item Erlaubt allen die Verwendung gleicher Programme, ohne Berücksichtigung
      von sozialem oder finanziellem Status, und fördert damit Zusammenarbeit
    \item Schützt vor sozialer Ausgrenzung
    \end{itemize}
  \end{block}

  \onslide<+->

  \begin{block}{Angepasste Kopien weitergeben zu können}
    \begin{itemize}
    \item Erlaubt allen Beteiligten, von Anpassungen anderer zu profitieren
    \item Schützt vor Abhängigkeiten bei Fehlerkorrekturen und vor unerwünschtem
      Verhalten der Software
    \end{itemize}
  \end{block}
\end{frame}

\begin{frame}
  \frametitle{Weiterführendes zu Freier Software}

  \onslide<+->

  \begin{block}{Mehr Informationen zu Freier Software}
    \begin{itemize}
    \item Free Software Foundation (\url{https://fsf/org})
    \item Free Software Foundation Europe (\url{https://fsfe.org})
    \item Ausführliche Beschreibung der vier Freiheiten:
      \url{https://www.gnu.org/philosophy/free-sw.html}
    \end{itemize}
  \end{block}

\end{frame}


\begin{frame}
  \frametitle{Freie Inhalte für die Schule}

  \onslide<+->

  \begin{block}{Probleme mit Urheberrecht}
    \begin{itemize}
    \item Fragen zum Urheberrecht in Schule relevant (Arbeitsblätter,
      Aufgabensammlungen, \dots)
    \item Oft unklar, ob Kopien für Unterricht genutzt werden dürfen
      % http://www.urheberrecht.de/schule/
    \end{itemize}
  \end{block}

  \onslide<+->

  \begin{block}{Freie Inhalte}
    \begin{itemize}
    \item Freiheiten für Software auf Unterrichtsmaterialien übertragen
    \item Lizenzen, die Weitergabe und Veränderungen ermöglichen \\
      (Creative Commons)
    \item Ergebnis: \emph{Open Educational Resources}
    \item Beispiele
      \begin{itemize}
      \item Zentrale für Unterrichtsmedien im Internet
        (\url{https://www.zum.de})
      \item Serlo (\url{https://de.serlo.org})
      \item \url{https://www.tutory.de/}
      \item Wikipedia
      \end{itemize}
    \end{itemize}
  \end{block}

\end{frame}

\begin{frame}

  \onslide<+->

  % stolen from Datenspuren lightning talk

  \begin{center}
    \url{https://fsfw-dresden.de/}
    $\;\;\left\{\;\;\text{
        \parbox{2.3cm}{
          \texttt{sprechstunde}\\[1mm]
          \texttt{mitmachen}\\[1mm]
          \texttt{fork}\\[1mm]
          \texttt{newsletter}
        }}
    \right.$

    \vspace*{2\bigskipamount}

    \includegraphics[width=50mm]{fsfw-netzwerke}
  \end{center}

\end{frame}

\end{document}

% Local Variables:
% TeX-engine: luatex
% End:

%  LocalWords:  Nutzungsvereinbarungen
